\section{Background shape: significance versus trial mass}
\label{sec:trial}

For one loose selection (DIRA$>0.995$, \Ls$>2$, \trialN{} candidates, 10\% skim) the fixed 15\,MeV Gaussian is moved in
5\,MeV steps from 1.74 to 1.99\,GeV/$c^2$ and fitted with pol2, pol3 and pol4 backgrounds in 1.66--2.08\,GeV/$c^2$
(Fig.~\ref{fig:trial}).

\begin{figure}[H]
\centering
\includegraphics[width=\textwidth]{trial_mass}
\caption{Left: $K\pi$ mass spectrum for DIRA$>0.995$, \Ls$>2$ (10\% skim, 5\,MeV bins; dotted line: PDG \Dzero{} mass).
Right: $z=S/\sigma_S$ of a 15\,MeV Gaussian fixed at the trial mass, for pol2 (black solid), pol3 (blue dashed) and pol4
(red dotted) background.}
\label{fig:trial}
\end{figure}

\paragraph{Result.} With pol3 the local significance runs smoothly from \trialZminThree{} to \trialZmaxThree{} (maximum
at \trialMmaxThree\,GeV/$c^2$); pol2 gives \trialZminTwo{} to \trialZmaxTwo, pol4 \trialZminFour{} to \trialZmaxFour.
At the \Dzero{} mass $z=\trialZTwo$ / \trialZThree{} / \trialZFour{} with $\chi^2/\mathrm{ndf}=\trialChiTwo$ /
\trialChiThree{} / \trialChiFour{} for pol2 / pol3 / pol4. A Gaussian of 15\,MeV width cannot produce a structure that
extends over about 200\,MeV (read from Fig.~\ref{fig:trial}); the large $|z|$ values measure the difference between the
spectrum and a low-order polynomial, not a peak. Polynomial-background fits on the raw spectrum therefore cannot establish
a \Dzero{} signal in loose selections; this motivates the ratio method of Sec.~\ref{sec:ratio}.
